\section{隐空间中的语义操作}

\subsection{确定性编码与解码}

DDIM 确定性采样建立的 $\mathbf{x}_T \leftrightarrow \mathbf{x}_0$ 双射映射，使得我们不仅可以从 $\mathbf{x}_T$ 生成 $\mathbf{x}_0$（解码），还可以从 $\mathbf{x}_0$ 恢复 $\mathbf{x}_T$（编码）。编码过程通过运行 DDIM 的前向方向实现：

$$
\mathbf{x}_{t+1} = \sqrt{\bar{\alpha}_{t+1}}\,\hat{\mathbf{x}}_0(\mathbf{x}_t) + \sqrt{1-\bar{\alpha}_{t+1}}\cdot\frac{\mathbf{x}_t - \sqrt{\bar{\alpha}_t}\,\hat{\mathbf{x}}_0(\mathbf{x}_t)}{\sqrt{1-\bar{\alpha}_t}}
$$

这意味着对于任意真实图片 $\mathbf{x}_0$，我们可以找到其在噪声空间中的``对应点'' $\mathbf{x}_T$，从该点出发再执行 DDIM 解码可以精确恢复原图。

\begin{importantbox}{DDIM 隐空间的语义结构}
由于确定性映射的存在，噪声空间 $\mathbf{x}_T$ 实际上成为了一个具有语义结构的隐空间。在这个空间中：
\begin{itemize}
    \item 相似的图像在隐空间中的表示也相近
    \item 隐空间中的线性运算（如插值）在数据空间中对应有意义的语义变化
    \item 不同于 VAE 的隐空间（低维），DDIM 的隐空间与数据空间维度相同
\end{itemize}
\end{importantbox}

\subsection{隐空间插值}

DDIM 论文展示了隐空间插值的有趣应用。给定两张真实图片 $\mathbf{x}_0^{(1)}$ 和 $\mathbf{x}_0^{(2)}$：

\begin{enumerate}
    \item 通过 DDIM 编码分别得到 $\mathbf{x}_T^{(1)}$ 和 $\mathbf{x}_T^{(2)}$
    \item 在隐空间中进行球面线性插值（SLERP）：
    $$
    \mathbf{x}_T^{(\lambda)} = \frac{\sin((1-\lambda)\theta)}{\sin\theta}\,\mathbf{x}_T^{(1)} + \frac{\sin(\lambda\theta)}{\sin\theta}\,\mathbf{x}_T^{(2)}, \quad \lambda \in [0, 1]
    $$
    其中 $\theta = \arccos\left(\frac{\langle\mathbf{x}_T^{(1)}, \mathbf{x}_T^{(2)}\rangle}{\|\mathbf{x}_T^{(1)}\|\,\|\mathbf{x}_T^{(2)}\|}\right)$
    \item 对插值点 $\mathbf{x}_T^{(\lambda)}$ 执行 DDIM 解码得到插值图片
\end{enumerate}

\begin{knowledgebox}{为什么使用球面插值（SLERP）而非线性插值}
由于隐空间中的噪声向量 $\mathbf{x}_T$ 是从标准高斯分布中采样的，在高维空间中，这些向量近似均匀分布在一个超球面上（范数集中在 $\sqrt{d}$ 附近）。线性插值（LERP）会产生范数较小的中间点，偏离高斯分布的典型集合（typical set），可能导致质量下降。球面插值（SLERP）保持向量范数不变，生成的中间点仍在超球面上，因此更合理。
\end{knowledgebox}

\subsection{与其他隐空间的对比}

\begin{table}[H]
\centering
\caption{不同模型隐空间特性对比}
\begin{tabular}{lccc}
\toprule
\textbf{模型} & \textbf{隐空间维度} & \textbf{编码确定性} & \textbf{插值质量} \\
\midrule
VAE & 低维 & 随机 & 中等 \\
GAN & 低维 & 通常无编码器 & 好 \\
DDIM & 与数据等维 & 确定性 & 好 \\
\bottomrule
\end{tabular}
\end{table}

\begin{warningbox}{DDIM 编码的精度依赖步数}
DDIM 编码过程本质上是 ODE 积分的数值近似。使用的步数越少，数值误差越大，``编码$\to$解码''的重建质量越低。在需要精确重建的应用（如图像编辑）中，通常需要使用较多的步数（如 100-200 步）来保证编码精度。
\end{warningbox}

\subsection{本章小结}

DDIM 的确定性采样构建了噪声空间到数据空间的双射映射，使噪声空间获得了语义结构。通过编码-插值-解码的流程，可以实现高质量的图像插值。球面线性插值（SLERP）比线性插值更适合高斯分布的隐空间。


